\documentclass[11pt]{article}
\usepackage[margin=1in]{geometry}
\usepackage[T1]{fontenc}
\usepackage[utf8]{inputenc}
\usepackage{amsmath,amssymb}
\usepackage{xcolor}
\usepackage{microtype}
\usepackage[colorlinks=true,linkcolor=blue!55!black,citecolor=blue!55!black,urlcolor=blue!55!black]{hyperref}

\newcommand{\B}{\mathbb B}
\newcommand{\N}{\mathbb N}

\title{Literature status: the two final mean-ergodicity questions in arXiv:2010.06938}
\author{Answers identified in arXiv:2206.00341}
\date{12 August 2026}

\begin{document}
\maketitle

\begin{center}
\fbox{\parbox{0.92\textwidth}{\raggedright\textbf{Status: answered in later literature.}
Questions~3.15 and 3.16 of arXiv:2010.06938 are both settled by
Theorems~1.3 and 1.4 of arXiv:2206.00341.  In particular, mean and uniform
mean ergodicity coincide for every composition operator on
$H^\infty(\B_n)$, and a symbol without an interior fixed point never induces
a mean-ergodic composition operator.}}
\end{center}

\section{Questions in the source}

Let $\varphi\colon\B_n\to\B_n$ be holomorphic, let
$C_\varphi f=f\circ\varphi$, and, in the interior-fixed-point case, let
$\rho$ be the holomorphic retraction associated with $\varphi$.
Theorem~3.6 of Keshavarzi~\cite{source} proves that if, for some $k\in\N$,
\[
 \sup_{z\in\B_n}\beta\bigl(\varphi^{kj}(z),\rho(z)\bigr)\longrightarrow0,
 \tag{B}
\]
then $C_\varphi$ is uniformly mean ergodic on $H^\infty(\B_n)$.
On PDF page~11, the paper then asks:
\begin{quote}
\textbf{Question 3.15.} Does the inverse [converse] of Theorem 3.6 hold?

\textbf{Question 3.16.} Is it true that composition operators induced by
holomorphic functions that have no interior fixed points are not mean
ergodic (extending Theorem 3.14)?
\end{quote}

\section{Exact later answers}

Keshavarzi and Hedayatian~\cite{answer} give the following characterization
as Theorem~1.3 on PDF page~2.  If $\varphi$ has an interior fixed point, then
the following are equivalent:
\begin{enumerate}
\item $C_\varphi$ is mean ergodic on $H^\infty(\B_n)$;
\item $C_\varphi$ is uniformly mean ergodic there;
\item for some $k\in\N$,
$\|\varphi^{kj}-\rho_\varphi\|_\infty\to0$.
\end{enumerate}
This also supplies the \emph{exact} converse requested in Question~3.15.
Indeed, uniform mean ergodicity implies mean ergodicity and hence condition
(3).  In the proof of $(3)\Rightarrow(2)$ (PDF pages~4--6), Theorem~1.2 and
the Kobayashi/Bergman estimate convert condition (3) into (B).  Together with
the source's Theorem~3.6, this yields
\[
 C_\varphi\text{ uniformly mean ergodic}
 \quad\Longleftrightarrow\quad
 \text{(B) holds for some }k.
\]

Theorem~1.4 of the later paper, also stated on PDF page~2 and proved on PDF
pages~10--11, says directly that if $\varphi$ has no interior fixed point,
then $C_\varphi$ is not mean ergodic on $H^\infty(\B_n)$.  The paragraph
immediately before the theorem explicitly identifies it as the answer to
the source's Question~3.16.

The proof addresses the fixed-witness obstruction missing from the source's
uniform-ergodicity argument.  It restricts the orbit operators to the first
coordinate copy of $H^\infty(\mathbb D)$.  Strong convergence there would,
using its Grothendieck--Dunford--Pettis structure and an operator-theoretic
bounded-below argument, force operator-norm convergence.  Boundary
Denjoy--Wolff convergence and the usual peak functions then contradict that
norm convergence.

\section{Scope and provenance}

Both concrete questions are answered in full.  The supporting authors knew
the provenance: their Theorem~1.4 explicitly cites Question~3.16, and their
introduction describes Theorem~1.3 as completing the interior-fixed-point
case left by the first author's earlier paper.  The broader side question
whether $H^\infty(\B_n)$ itself is a Grothendieck space with the
Dunford--Pettis property is not answered by these theorems and is not claimed
here.  This packet records a literature resolution, not a new proof.

\begin{thebibliography}{2}
\bibitem{source}
H.~Keshavarzi,
\emph{Mean ergodic composition operators on $H^\infty(\mathbb B_n)$},
arXiv:2010.06938; Positivity 26 (2022), article 30,
\href{https://doi.org/10.1007/s11117-022-00901-5}{doi:10.1007/s11117-022-00901-5}.

\bibitem{answer}
H.~Keshavarzi and K.~Hedayatian,
\emph{Ergodic behaviors of composition operators acting on space of bounded
holomorphic functions}, arXiv:2206.00341v1, 2022.
\end{thebibliography}

\end{document}
